\documentclass{article}

\usepackage{neurips_2026}
\usepackage[utf8]{inputenc}
\usepackage[T1]{fontenc}
\usepackage{amsmath,amssymb}
\usepackage{booktabs}
\usepackage{microtype}
\usepackage{array}
\usepackage{tabularx}
\usepackage{url}
\usepackage{hyperref}

\title{When Does GRPO Have a Learning Signal?\\
A Claim-to-Evidence Audit of Group-Relative RL}

\author{}

\newcommand{\zvf}{\mathrm{ZVF}}
\newcommand{\gu}{\mathrm{GU}}
\newcommand{\usable}{P_{\mathrm{usable}}}
\newcolumntype{Y}{>{\raggedright\arraybackslash}X}

\begin{document}
\maketitle

\begin{abstract}
We study when a critic-free group-relative RL loop has a usable learning
signal. Across low-budget tool, code, browser-control, and math runs, the
common mechanism is within-group reward contrast: all-wrong groups produce
cold-start collapse, while all-correct groups produce saturation. We
operationalize this with Zero-Variance Fraction (ZVF) and Gradient Utilization
(GU) as early triage diagnostics, not calibrated predictors. The clean
held-out GSM8K control shows only a small, non-significant improvement over the
base model, and paired disagreement tests on the same 200 examples do not show
a reliable improvement for any seed. Tool, code, and browser-control results
are custom, subset-bound, or integration-only proxy measurements. The
contribution is therefore not a GRPO leaderboard result, but a
claim-to-evidence audit showing how reward density, base-policy hit rate, group
size, and stack details determine whether a nominal GRPO run has interpretable
signal.
\end{abstract}

\section{Introduction}

This paper does not claim that Group Relative Policy Optimization (GRPO)
broadly improves reasoning. The narrower claim is that in a critic-free
group-relative update, the learning signal exists only when sampled groups
contain reward diversity. If every completion for a prompt is wrong, the group
has no ordering signal; if every completion is right, the group is saturated.
Both cases can produce flat or misleading reward curves. We audit this failure
mode in a heterogeneous post-training corpus covering mathematical reasoning,
tool-call formatting, code tasks, and a BrowserGym/MiniWoB integration smoke
test.

The paper is organized as a claim-to-evidence audit rather than a benchmark.
Online training reward is interpreted as evidence about optimization dynamics,
not as benchmark accuracy. Held-out evaluations are used for capability claims.
Tool-use rewards are treated as schema or proxy rewards unless tools are
executed and evaluated end to end. Algorithm labels such as PPO, GRPO, and DPO
are treated as incomplete treatment descriptions unless the full stack is
reported.

\begin{table}[t]
\centering
\caption{Claim/evidence/scope hierarchy for the compact submission.}
\label{tab:claim_scope}
\small
\begin{tabularx}{\linewidth}{p{0.22\linewidth}Y Y Y}
\toprule
Claim & Evidence & Allowed conclusion & Not allowed \\
\midrule
Reward contrast governs signal &
Mixed-group probability, early ZVF/GU, collapse and saturation traces &
Group-relative updates need within-prompt reward variation &
Universal GRPO improvement \\
Training reward is not capability &
GSM8K base 82.0\% vs. GRPO mean 83.3\%, $p=0.26$; paired seed tests below &
Held-out math gain is not established &
``GRPO improves GSM8K'' \\
Stack labels are incomplete &
Backend, parser, sampler, loss-mask, LoRA, KL/reference, checkpoint, and evaluator differences &
Report stack-conditioned results &
PPO/GRPO/DPO leaderboard ranking \\
Tool and browser results are proxy evidence &
JSON/tool-name/argument-key rewards; BrowserGym smoke reward 0.0 &
Schema/integration evidence only &
End-to-end tool or browser competence \\
\bottomrule
\end{tabularx}
\end{table}

Our contributions are: (1) a reward-diversity diagnostic for group-relative
RL using ZVF and GU; (2) a finite-group mechanism that connects base-policy hit
rate, group size, reward density, and task difficulty; (3) a negative held-out
capability result on GSM8K; and (4) a stack-conditioned reporting template for
post-training comparisons.

\section{Mechanism: Reward Contrast}

For a prompt $x$, let $p_x$ be the probability that the current policy samples
a rewarded completion under the rollout distribution, and let $G$ be group
size. With a binary verifier, a sampled group has usable contrast if it
contains at least one success and at least one failure. The probability of such
a group is
\begin{equation}
  \usable(x;G) = 1 - (1-p_x)^G - p_x^G .
\end{equation}
For a prompt set, the empirical design quantity is
\begin{equation}
  \widehat{\usable}(G) =
  \frac{1}{N}\sum_{x=1}^N \left[1-(1-p_x)^G-p_x^G\right].
\end{equation}
This is not a theorem that GRPO improves capability. It is a design rule for
whether the sampled rollout distribution can produce a non-zero relative
advantage.

We log two diagnostics. Zero-Variance Fraction is the fraction of prompt groups
whose rewards are all identical:
\[
\zvf = \Pr(\operatorname{Var}(r_{x,1:G})=0).
\]
Gradient Utilization is the fraction of rollout groups or steps that yield a
non-zero advantage:
\[
\gu = 1-\zvf .
\]
High ZVF with low reward indicates cold-start collapse; high ZVF with high
reward indicates saturation. Low ZVF indicates that rollout compute is
producing mixed groups. We use ZVF/GU for triage, not for calibrated
forecasting: pooled lagged analyses in the artifact do not establish that ZVF
adds reliable predictive power beyond current reward.

\section{Audit Corpus and Evidence Levels}

The reviewer-facing paper uses a compact ledger. It does not treat run count
as a contribution. The broader repository contains interrupted attempts,
partial probes, local W\&B/Hugging Face/Tinker mirrors, and bookkeeping
records; those are preservation artifacts, not independent experiments. The
main paper uses only the evidence classes in Table~\ref{tab:evidence_levels}.

\begin{table}[t]
\centering
\caption{Evidence levels used in the main text.}
\label{tab:evidence_levels}
\small
\begin{tabularx}{\linewidth}{p{0.23\linewidth}Y Y}
\toprule
Evidence type & Examples in this paper & Supports \\
\midrule
Held-out capability &
Base and GRPO checkpoints on the same 200 GSM8K test examples &
Capability or generalization claims, when paired or controlled \\
Online training reward &
GSM8K/MATH/tool/code rollout rewards during training &
Optimization dynamics and reward-environment behavior \\
Proxy/schema reward &
JSON validity, tool-name match, argument-key presence &
Structured-output behavior under a custom parser \\
Integration smoke test &
BrowserGym/MiniWoB routing through browser, W\&B, and checkpoint plumbing &
Systems integration only; no learned browser success \\
Stack case study &
Runs that differ in backend, parser, sampler, LoRA, loss, or evaluator &
Treatment-identifiability concerns, not algorithm ranking \\
\bottomrule
\end{tabularx}
\end{table}

The group-relative runner is GRPO-inspired, not a canonical reproduction of
DeepSeekMath GRPO. Some runs omit reference KL, PPO clipping, or completion-only
masking, and the managed backend hides some low-level implementation details.
Those are part of the treatment and therefore part of the claim boundary.

\section{Results}

\subsection{Reward Diversity Explains Signal Availability}

The strongest empirical pattern is diagnostic. Strict sparse-reward tool runs
without warm-starting produce all-zero reward groups and no useful update.
High-hit-rate math settings rapidly saturate, producing all-correct groups.
Intermediate regimes produce mixed groups, and larger group sizes preserve
usable contrast longer.

\begin{table}[t]
\centering
\caption{Group size changes online reward dynamics for one 8B GSM8K setup.
These are training-rollout rewards, not held-out benchmark accuracies.}
\label{tab:group_size}
\small
\begin{tabular}{lcccc}
\toprule
$G$ & First-5 reward & Last-10 reward & Peak reward & Zero-reward steps \\
\midrule
4  & 32.5\% & 23.8\% & 62.5\%  & 20\% \\
8  & 33.8\% & 24.4\% & 68.8\%  & 6\% \\
16 & 29.4\% & 36.2\% & 75.0\%  & 2\% \\
32 & 32.8\% & 54.7\% & 100.0\% & 0\% \\
\bottomrule
\end{tabular}
\end{table}

The same table should not be read as a capability result: it is a
training-environment measurement under the rollout parser. Its value is that it
shows the finite-group mechanism in action. Increasing $G$ raises the chance
that a prompt group contains at least one success and one failure; once all
groups are all-correct or all-wrong, the relative advantage loses ordering
information.

\subsection{Held-Out GSM8K Does Not Establish Capability Gain}

The cleanest capability check is a 200-example GSM8K held-out slice evaluated
with identical greedy decoding for the base model and five GRPO checkpoints.
The base model solves 164/200 examples (82.0\%). GRPO checkpoints average
83.3\% with standard deviation 2.2 percentage points. A one-sample test against
the base accuracy gives $p=0.26$, so the held-out lift is not statistically
reliable.

Because the base and GRPO checkpoints are evaluated on the same examples, a
paired disagreement analysis is more informative than a headline mean. It asks
how many examples are solved only by the base model versus only by the GRPO
checkpoint.

\begin{table}[t]
\centering
\caption{Paired GSM8K disagreement audit on the same 200 held-out examples.
Exact paired tests show example churn, not reliable improvement.}
\label{tab:paired_gsm8k}
\small
\begin{tabular}{lccc}
\toprule
GRPO seed & Base-only correct & GRPO-only correct & Exact paired $p$ \\
\midrule
42  & 13 & 15 & 0.851 \\
137 & 15 & 16 & 1.000 \\
256 & 19 & 16 & 0.736 \\
512 & 12 & 16 & 0.572 \\
999 & 7  & 16 & 0.093 \\
\bottomrule
\end{tabular}
\end{table}

The conclusion is deliberately negative: GRPO changes which examples are
solved, but this held-out slice does not support a broad math-capability gain.
This negative result is the credibility anchor for the rest of the audit.

\subsection{Tool, Code, and Browser Evidence Is Proxy-Bound}

Tool-call training can improve structured output under a custom reward, but the
positive rows measure parser-aligned behavior rather than executed tool
competence. In one SFT-warmed tool pipeline, JSON validity rose from 0\% to
92\%, correct tool-name reward reached 50\%, and argument-key presence reached
42\%. These are useful schema-compliance signals, but not a standardized
ToolBench/Gorilla-style held-out API evaluation.

HumanEval evidence is also limited. A 50-problem subset moved from 16/50 to
20/50 pass@1, but the Fisher exact test is non-significant ($p=0.53$), and the
full canonical HumanEval pass@$k$ harness was not completed in the central
setup. The BrowserGym/MiniWoB run is even narrower: it verified browser
observation/action/screenshot plumbing through the training pipeline, but a
one-step smoke run received 0.0 reward and no trainable datum groups.

\begin{table}[t]
\centering
\caption{Agentic evidence levels. The paper includes browser control as
systems integration evidence, not learned browser-task competence.}
\label{tab:agentic}
\small
\begin{tabularx}{\linewidth}{p{0.24\linewidth}Y Y}
\toprule
Level & Evidence & Interpretation \\
\midrule
Schema validity & SFT-warmed JSON/tool-call rewards & Structured-output behavior can be shaped after warm-starting \\
Tool choice proxy & Tool-name and argument-key reward & Parser-aligned proxy, not full execution correctness \\
Browser routing & BrowserGym/MiniWoB smoke & Browser pipeline integrated; reward remained 0.0 \\
Environment success & Not established & Requires standardized BrowserGym/WebArena suites and baselines \\
\bottomrule
\end{tabularx}
\end{table}

\subsection{Algorithm Labels Are Not Treatment Definitions}

The artifact contains PPO/GRPO/DPO comparisons, but they are not used as
algorithm rankings. In LLM post-training, the treatment is the tuple
\[
T=(A,M,C,R,S,O,L,K,E),
\]
where $A$ is the nominal algorithm, $M$ the checkpoint/tokenizer, $C$ the prompt
construction, $R$ the reward/parser, $S$ the sampler and group construction,
$O$ the optimizer schedule, $L$ the loss mask and KL/reference handling, $K$
the backend/kernel/precision stack, and $E$ the evaluator and checkpoint rule.
Two runs with the same $A=\mathrm{GRPO}$ but different tuple elements are not
replicates of one treatment. The paper therefore claims treatment
under-identification, not framework superiority.

\section{Related Work}

DeepSeekMath introduced GRPO as a critic-free, group-relative alternative to
PPO for mathematical reasoning \citep{shao2024deepseekmath,schulman2017proximal}.
Our paper does not reproduce that recipe; it audits a GRPO-inspired runner and
asks when its sampled groups contain usable reward contrast. Zero-variance
prompt work and online difficulty filtering are the closest algorithmic
relatives: they show that all-correct or all-incorrect groups are poor learning
substrates and that intermediate-difficulty prompts carry more signal
\citep{rlzvp2025,bae2025online}. Spurious-reward results motivate our evidence
hierarchy by showing that reward improvements can amplify model-specific priors
without establishing general capability \citep{shao2025spurious}.

The stack-identifiability claim extends older deep-RL reproducibility findings:
seed noise, implementation details, and statistical uncertainty can dominate
nominal algorithm labels \citep{henderson2018deep,agarwal2021deep}. Tool-use
benchmarks such as Gorilla and ToolLLM evaluate held-out API behavior
\citep{patil2023gorilla,qin2024toolllm}; ToolRM/FC-RewardBench studies reward
models for tool calling \citep{wang2025toolrm}. Our custom tool rewards are not
equivalent to those standardized evaluations.

\section{Limitations}

\paragraph{Noncanonical runner.}
The experiments use a GRPO-inspired group-relative policy-gradient runner. They
should not be read as claims about canonical DeepSeekMath GRPO or every TRL,
verl, OpenRLHF, or Tinker implementation.

\paragraph{Short horizons and limited replication.}
Several results are single-run, short-budget, or interrupted. Group-size and
architecture observations are design probes. They motivate the reward-contrast
mechanism but do not establish universal scaling laws or capacity thresholds.

\paragraph{Proxy rewards.}
Tool-use scores are parser/schema rewards unless execution-level tool success
is measured. BrowserGym/MiniWoB evidence is integration-only. HumanEval evidence
is subset-bound. Online reward curves are dynamics evidence, not capability.

\paragraph{Public benchmark contamination and missing baselines.}
GSM8K and HumanEval are public and may appear in pretraining data. The clean
GSM8K comparison is therefore a narrow same-slice control, not proof of novel
reasoning. We do not include full matched PPO, REINFORCE, DPO, grammar-constrained
decoding, or supervised-only baselines for every setting.

\paragraph{Diagnostics, not predictors.}
ZVF/GU are useful alarms for cold-start collapse and saturation. They are not
claimed to be calibrated predictors of final capability, and the artifact's
lagged analysis does not justify that stronger claim.

\section{Conclusion}

The strongest conclusion is methodological: GRPO-style critic-free RL should be
audited for reward contrast before reward curves are interpreted. In this
corpus, rising online reward sometimes reflects schema compliance, parser
alignment, or saturation, while the clean held-out GSM8K control remains
non-significant. The practical contribution is a reporting discipline: log
ZVF/GU, separate training reward from held-out capability, and define the full
post-training stack before making algorithm-level comparisons.

\bibliographystyle{plainnat}
\bibliography{references}

\clearpage
\section*{NeurIPS Paper Checklist}
\begin{enumerate}
\item \textbf{Claims.} The abstract and introduction state that this is a
diagnostic audit, not a broad GRPO capability claim.
\item \textbf{Limitations.} Section 6 lists runner scope, short horizons,
proxy rewards, public benchmark contamination, missing baselines, and diagnostic
scope.
\item \textbf{Reproducibility.} The artifact includes code, result ledgers, and
checksums. Managed-backend details that cannot be fully reproduced are treated
as part of the treatment boundary.
\item \textbf{Statistics.} The main held-out claim reports both the
non-significant aggregate GSM8K comparison and paired disagreement tests.
\item \textbf{Datasets.} GSM8K and HumanEval are public benchmarks; tool and
browser tasks are custom/proxy setups and are labeled as such.
\item \textbf{Ethics.} The paper is an audit of post-training methodology. It
does not release a deployed tool-use or browser-control system and does not
claim end-to-end autonomous competence.
\end{enumerate}

\end{document}
